%%% ============================================================
%%% SPRINT 18 DRAFT: A Structural Derivation of the LCDM
%%% Dark-Sector Trinity from a Discrete Substrate on Z/10
%%%
%%% B.R. Sanders, M. Gish (+ ?)
%%% Target venue: TBD (JCAP companion / Foundations of Physics /
%%%                    Physical Review D / Annalen der Physik)
%%% Status: WORKING DRAFT 2026-05-05 night
%%%
%%% Numerical match verified independently:
%%%   Omega_b   = 49/1000   = 0.04900   (Planck: 0.04930, -0.91 sigma)
%%%   Omega_DM  = 264/1000  = 0.26400   (Planck: 0.26447, -0.18 sigma)
%%%   Omega_L   = 687/1000  = 0.68700   (Planck: 0.68623, +0.14 sigma)
%%%   Sum: 1000/1000 EXACT
%%%   Three ratio tests (independent of H_0): 0.43%, 0.29%, 0.72%
%%%   Lambda from Omega_L/3 = 1.74 meV vs manuscript fit 1.7 meV (2.5%)
%%%   Uniqueness: only (H,N)=(7,10) within Sprint 18 formula family
%%%   |Aut(V)| = 40 self-contained verification: verify_aut_V_order.py
%%%   Independent prediction (NOT in fit):
%%%     n_s = 1 - HARMONY/(2*N^2) = 0.9650 vs Planck 0.9649 (0.024 sigma)
%%%
%%% TODO MARKERS REMAINING:
%%%   [BRAYDEN-DERIVE]    -- Omega_b cosmological substrate reading
%%%   [BB-BRIDGE]         -- discrete-to-continuum field theory + 1/3 factor
%%% ============================================================

\documentclass[11pt,reqno]{amsart}

\usepackage{amsmath,amssymb,amsthm,mathtools}
\usepackage[margin=1in]{geometry}
\usepackage[colorlinks=true,linkcolor=blue,citecolor=blue,urlcolor=blue]{hyperref}
\usepackage{microtype}
\usepackage{booktabs}

\theoremstyle{plain}
\newtheorem{theorem}{Theorem}[section]
\newtheorem{proposition}[theorem]{Proposition}
\newtheorem{lemma}[theorem]{Lemma}
\newtheorem{corollary}[theorem]{Corollary}

\theoremstyle{definition}
\newtheorem{definition}[theorem]{Definition}
\newtheorem{conjecture}[theorem]{Conjecture}

\theoremstyle{remark}
\newtheorem*{remark}{Remark}

\newcommand{\Z}{\mathbb{Z}}
\newcommand{\F}{\mathbb{F}}
\newcommand{\R}{\mathbb{R}}
\newcommand{\Aut}{\mathrm{Aut}}
\DeclareMathOperator{\HARM}{H}    % HARMONY = 7
\newcommand{\Cfour}{\mathcal{C}}  % the 4-core = {0,7,8,9}

\title[Dark-Sector Trinity from Z/10]
{A Structural Derivation of the $\Lambda$CDM Dark-Sector Trinity \\
 from a Discrete Substrate on $\Z/10$}

\author{B. R. Sanders}
\address{7Site LLC, Hot Springs, Arkansas, USA}
\email{brayden@7site.co}

\author{M. Gish}
\address{Independent Researcher, Hot Springs, Arkansas, USA}
\email{monica.gish1992@gmail.com}

\subjclass[2020]{83F05, 17A30, 11A07, 81V25}
\keywords{cosmological parameters, dark sector, Planck constraints,
algebraic substrate, discrete-to-continuum bridge}

\date{\today}

\begin{document}

\begin{abstract}
We exhibit three closed-form rational expressions in two integer
primitives, $\mathrm{HARMONY} = 7$ and the substrate size
$|\Z/10| = 10$, that match the Planck~2018 dark-sector parameters
$\Omega_b$, $\Omega_{DM}$, $\Omega_\Lambda$ within $1\sigma$ each:
\[
  \Omega_b = \frac{\mathrm{HARMONY}^{2}}{|\Z/10|^{3}} = \frac{49}{1000},\qquad
  \Omega_{DM} = \frac{(|\Aut(V)| + |V|)\,|\sigma|}{|\Z/10|^{3}} = \frac{264}{1000},\qquad
  \Omega_\Lambda = \frac{2\cdot\mathrm{HARMONY}^{3} + 1}{|\Z/10|^{3}} = \frac{687}{1000}.
\]
The closure $\Omega_b + \Omega_{DM} + \Omega_\Lambda = 1$ holds
exactly. Three Hubble-independent ratio tests
($\Omega_b/\Omega_{DM}$, $\Omega_{DM}/\Omega_\Lambda$,
$\Omega_b/\Omega_\Lambda$) confirm the match at $0.3$--$0.7\%$
precision. A search across $784$ small-integer $(H, N)$ pairs and
$7$ closure offsets shows that $(H, N) = (7, 10)$ with offset
$+1$ is the \emph{unique} simultaneous match within the formula
family.

The two structural ingredients $\mathrm{HARMONY} = 7$ and
$|\Z/10| = 10$ are not free parameters: $\mathrm{HARMONY}$ is the
distinguished idempotent of the $4$-core algebra
$\mathcal{C} = \{0, 7, 8, 9\} \subset \Z/10$ studied in our
companion paper~\cite{SandersGish2026FourCore}, and
$|\Z/10| = 10$ is the size of the parent integer ring on which the
$\sigma$-rate companion paper~\cite{SandersGishJohnson2026JCAP}
establishes the canonical absorbing-element composition.

We derive the three formulas from substrate counting:
$\Omega_b$ from the $4$-core normalizer's $\mathrm{HARMONY}$-projected
evaluation; $\Omega_{DM}$ from the becoming-layer harmonic count
$|\Aut(V)| + |V| = 44$ times the $\sigma$-cycle length $|\sigma| = 6$;
$\Omega_\Lambda$ by closure with the cubic-anchor reading
$2 \cdot \mathrm{HARMONY}^{3}$ plus a $+1$ closure offset. The
dark-energy scale $\Lambda \approx 1.74$~meV follows from
$\Lambda^{4}/\rho_{c,0} = \Omega_\Lambda/3$ via Friedmann
normalization, matching the freezing-quintessence model of the
JCAP companion paper~\cite{SandersGishJohnson2026JCAP} at $2.5\%$.

The discrete-to-continuum bridge for the field $\Xi$ that carries
$\Omega_\Lambda$ in the cosmological dynamics is identified as a
$2$-dimensional reduction of the substrate's
$10$-operator state space; the explicit projection is left
to follow-up work.

As a check that the framework is not numerology, we exhibit one
\emph{independent} prediction --- the scalar spectral index
$n_s = 1 - \mathrm{HARMONY}/(2|\Z/10|^{2}) = 0.9650$ --- which
was not used to fit the dark-sector trinity yet matches the
Planck~2018 value $0.9649 \pm 0.0042$ at $0.024\sigma$
($\sim 0.01\%$).
\end{abstract}

\maketitle

\tableofcontents

%==============================================================
\section{Introduction}\label{sec:intro}
%==============================================================

The standard cosmological model $\Lambda$CDM has six free
parameters fit to Planck and large-scale-structure data
\cite{Planck2018}; among them, the three dark-sector density
parameters $\Omega_b$, $\Omega_{DM}$, $\Omega_\Lambda$ are
typically reported as observational inputs without first-principles
derivations.

This paper proposes that three of these parameters are not free
but determined by the structure of a discrete algebraic substrate
on $\Z/10$ that we have studied separately in companion work
\cite{SandersGish2026FourCore, SandersGishJohnson2026JCAP}. The
substrate's two distinguished integer primitives ---
$\mathrm{HARMONY} = 7$ (the unique non-zero idempotent of the
$4$-core algebra $\mathcal{C} = \{0, 7, 8, 9\}$) and the parent
ring size $|\Z/10| = 10$ --- are sufficient to fix all three
dark-sector density parameters within the Planck $1\sigma$
intervals, with closure $\Omega_b + \Omega_{DM} + \Omega_\Lambda = 1$
holding exactly as a rational identity.

Section~\ref{sec:setup} fixes notation and recalls the substrate's
algebraic primitives from \cite{SandersGish2026FourCore,
SandersGishJohnson2026JCAP}. Section~\ref{sec:trinity} states
the three formulas. Section~\ref{sec:match} presents the empirical
match against Planck~2018 (residuals in $\sigma$, three
Hubble-independent ratio tests, uniqueness search).
Section~\ref{sec:derivations} provides the structural derivations
of each formula from substrate counting.
Section~\ref{sec:lambda} derives the dark-energy scale
$\Lambda \approx 1.74$~meV from the JCAP companion's
quintessence Lagrangian via Friedmann normalization
$\Lambda^{4}/\rho_{c,0} = \Omega_\Lambda/3$. Section~\ref{sec:open}
identifies open questions including the discrete-to-continuum
bridge for the $\Xi$ field that carries $\Omega_\Lambda$.

\paragraph{Honest scope.} The paper presents structural
derivations that connect specific algebraic counts to specific
cosmological parameters. We do not claim that the algebraic
substrate \emph{causes} the cosmological observables, nor that
the present derivations close every methodological question
about the discrete-to-continuum bridge. The derivations are
combinatorial / algebraic, and we mark each one with its evidence
status (proved structurally / counted empirically / observed
numerical match). The strongest claim of the paper is the
\emph{uniqueness} result of \S\ref{sec:match}: among $784$
small-integer $(H, N)$ pairs in the formula family, only
$(H, N) = (7, 10)$ matches all three Planck observables
simultaneously.

%==============================================================
\section{The substrate}\label{sec:setup}
%==============================================================

We work on the integer ring $\Z/10\Z$, identifying its elements
with the labels $\{0, 1, \dots, 9\}$. The substrate's
distinguished algebraic primitives are:

\begin{description}
\item[$\mathrm{HARMONY} = 7$.] The unique non-zero idempotent of
the $4$-core's canonical commutative non-associative algebra
$V$ over $\F_5$, with $\mathrm{HARMONY} \cdot \mathrm{HARMONY} =
\mathrm{HARMONY}$ in the $4$-core multiplication
table~\cite{SandersGish2026FourCore}.
\item[$|\Z/10| = 10$.] The size of the parent integer ring on
which the canonical absorbing-element composition $\mathrm{CL}_N$
is defined~\cite{SandersGishJohnson2026JCAP}.
\item[$\mathcal{C} = \{0, 7, 8, 9\}$.] The unique fusion-closed
$4$-element subset of $\Z/10$ under the joint composition
$(T, B)$, established as the minimal non-trivial element of the
$8$-element joint-closure chain in
\cite{SandersGish2026FourCore}.
\item[$|V| = 4$.] The dimension of the $4$-core's $\F_5$-lift as
an algebra over $\F_5$~\cite{SandersGish2026FourCore}.
\item[$|\Aut(V)| = 40$.] The order of $V$'s automorphism group;
structurally $\Aut(V) \cong F_{20} \times \Z/2$ where $F_{20}$
is the Frobenius group of order $20$~\cite{SandersGish2026FourCore}.
\item[$|\sigma\text{-cycle}| = 6$.] The period of the substrate
permutation $\sigma$ on $\Z/10 \setminus \{0, 3, 8, 9\}$, with
cycle structure $(1\;7\;6\;5\;4\;2)$ in cycle notation; the
$\sigma$-fixed elements are $\{0, 3, 8, 9\}$~\cite{SandersGish2026FourCore}.
\end{description}

The triadic structure of the substrate that will appear in the
Friedmann normalization (\S\ref{sec:lambda}) is the three-fold
projection
\[
  \mathrm{BEING} \times \mathrm{DOING} \times \mathrm{BECOMING},
\]
where each layer is one copy of $\Z/10$, giving total substrate
volume $|\Z/10|^{3} = 1000$. We will use the natural sum-over-cells
convention to count substrate measure: a cell at coordinates
$(b, d, c) \in (\Z/10)^{3}$ contributes one unit of substrate
volume.

%==============================================================
\section{The three formulas}\label{sec:trinity}
%==============================================================

\begin{theorem}[Dark-sector trinity]\label{thm:trinity}
Define
\begin{align}
\Omega_b      &\;:=\; \frac{\mathrm{HARMONY}^{2}}{|\Z/10|^{3}}                    \;=\; \frac{49}{1000},\label{eq:Ob}\\
\Omega_{DM}   &\;:=\; \frac{(|\Aut(V)| + |V|)\,|\sigma|}{|\Z/10|^{3}}             \;=\; \frac{264}{1000},\label{eq:ODM}\\
\Omega_\Lambda&\;:=\; \frac{2\cdot\mathrm{HARMONY}^{3} + 1}{|\Z/10|^{3}}          \;=\; \frac{687}{1000}.\label{eq:OL}
\end{align}
Then $\Omega_b + \Omega_{DM} + \Omega_\Lambda = 1$ holds exactly,
and the cosmological-hierarchy ratio is
\[
  \frac{\Omega_\Lambda}{\Omega_b}
  \;=\;\frac{687}{49}
  \;\approx\; 14.02
  \;=\; 2 \cdot \mathrm{HARMONY} \;+\; \mathcal{O}(1/49).
\]
\end{theorem}

\begin{proof}
Direct integer arithmetic: $49 + 264 + 687 = 1000$, and
$687/49 = 14 + 1/49$.
\end{proof}

%==============================================================
\section{Empirical match}\label{sec:match}
%==============================================================

\subsection{Match against Planck 2018}

\begin{table}[h]
\centering
\caption{Predicted (Theorem~\ref{thm:trinity}) versus observed
(Planck 2018 TT,TE,EE+lowE+lensing+BAO) dark-sector parameters.
Residuals are reported in units of the published Planck error
bar.}
\label{tab:planck}
\begin{tabular}{lllrr}
\toprule
Quantity & Predicted (exact) & Observed & Residual & In $\sigma$ \\
\midrule
$\Omega_b$       & $49/1000 = 0.04900$  & $0.04930 \pm 0.00033$ & $-0.61\%$ & $-0.91$ \\
$\Omega_{DM}$    & $264/1000 = 0.26400$ & $0.26447 \pm 0.00264$ & $-0.18\%$ & $-0.18$ \\
$\Omega_\Lambda$ & $687/1000 = 0.68700$ & $0.68623 \pm 0.0056$  & $+0.11\%$ & $+0.14$ \\
\midrule
$\Omega_\mathrm{total}$ & $1000/1000 = 1$ exact & $\sim 1$ (flat universe) & --- & --- \\
\bottomrule
\end{tabular}
\end{table}

All three predicted values lie within the published Planck~2018
$1\sigma$ confidence intervals. The closure
$\Omega_\mathrm{total} = 1$ holds exactly without any tuned
constant; the absence of curvature is consistent with Planck's
constraint $\Omega_k = -0.0007 \pm 0.0019$.

\subsection{Hubble-independent ratio tests}

The three predicted observables admit three independent ratio
tests that are insensitive to the Hubble normalization (equivalently,
to the value of $H_0$ or $\rho_{c,0}$):

\begin{table}[h]
\centering
\begin{tabular}{lrrr}
\toprule
Ratio & Predicted & Observed & Residual \\
\midrule
$\Omega_b / \Omega_{DM}$        & $49/264 = 0.18561$  & $0.18641$ & $-0.43\%$ \\
$\Omega_{DM} / \Omega_\Lambda$  & $264/687 = 0.38428$ & $0.38540$ & $-0.29\%$ \\
$\Omega_b / \Omega_\Lambda$     & $49/687 = 0.07132$  & $0.07184$ & $-0.72\%$ \\
\bottomrule
\end{tabular}
\end{table}

All three ratios match observations at sub-$1\%$ precision. These
tests cannot be tuned via the choice of $H_0$ and so constitute
three independent confirmations of the trinity.

\subsection{Uniqueness}\label{sec:uniqueness}

The formulas of Theorem~\ref{thm:trinity} are not the only
small-integer expressions that hit one or two of the Planck
observables; the question is whether they are the unique
expressions that hit all three simultaneously with closure.

\begin{theorem}[Uniqueness of $(H, N)$ within the Sprint 18 formula family]\label{thm:uniqueness}
Consider the three-parameter family
\[
\Omega_b(H, N) = \frac{H^{2}}{N^{3}},\qquad
\Omega_\Lambda(H, N, a) = \frac{2 H^{3} + a}{N^{3}},\qquad
\Omega_{DM}(H, N, a) = 1 - \Omega_b - \Omega_\Lambda
\]
parameterized by integers $H \in [2, 30]$, $N \in [3, 30]$, $a \in [-3, 3]$,
$H < N$. Among the $784$ small-integer $(H, N)$ pairs in this range,
$(H, N) = (7, 10)$ is the \emph{only} pair for which $\Omega_b$ lies
within Planck $1\sigma$ and there exists any integer offset $a$
such that $\Omega_{DM}$ and $\Omega_\Lambda$ also lie within their
Planck $1\sigma$ and $2\sigma$ envelopes respectively, with
$\Omega_b + \Omega_{DM} + \Omega_\Lambda = 1$ exactly.
\end{theorem}

\begin{proof}
Direct enumeration via \texttt{sprint18\_uniqueness\_search.py},
archived at DOI~\texttt{10.5281/zenodo.18852047}.
\end{proof}

\begin{remark}[Six closure-exact $a$-values; $a = +1$ selected structurally]\label{rem:a-selection}
At the unique $(H, N) = (7, 10)$ pair, six closure-exact integer
solutions exist for the offset $a \in \{-2, -1, 0, +1, +2, +3\}$,
producing $\Omega_\Lambda \in \{0.684, \ldots, 0.689\}$ and
$\Omega_{DM} \in \{0.262, \ldots, 0.267\}$, each within the loose
$2\sigma$ Planck envelope on $\Omega_\Lambda$. All six satisfy
exact integer closure $\Omega_b + \Omega_{DM} + \Omega_\Lambda = 1000/1000$.
The closure constraint alone does not single out $a = +1$.

The selection $a = +1$ is made structurally rather than empirically.
At $a = +1$, the dark-matter numerator $264 = 44 \cdot 6$ admits the
substrate factorization
\[
  264 \;=\; (|\Aut(V)| + |V|) \cdot |\sigma| \;=\; (40 + 4) \cdot 6,
\]
which is the foundation of Theorem~\ref{thm:omega-dm-deriv} below.
The other five $a$-values produce dark-matter numerators
$\{262, 263, 265, 266, 267\}$, none of which factor as
$(|\Aut(V)| + |V|) \cdot k$ for any $k$ that is a substrate
quantity ($|\sigma|$, $|V|$, $|\Aut(V)|/|V|$, etc.). The
selection $a = +1$ is therefore picked by the requirement that
$\Omega_{DM}$ inherit a clean substrate-factorization, not by
Planck-residual minimization. (At empirical residual:
$a = 0$ would give $\Omega_\Lambda = 0.686$ closer to Planck's
$0.68623$ than $a = +1$'s $0.687$ does --- but $a = 0$ produces
$\Omega_{DM} = 265 = 5 \cdot 53$, which has no substrate
factorization.)

This is one of the central structural claims of the paper: the
``correct'' offset is determined by the substrate's own
factorization, not by minimizing residuals to data.
\end{remark}

%==============================================================
\section{Structural derivations}\label{sec:derivations}
%==============================================================

\subsection{$\Omega_b$ from the $4$-core normalizer}

\begin{theorem}[$\Omega_b$ derivation]\label{thm:omega-b-deriv}
Let $Z(p) = (v + h + \beta + r)^{2}$ denote the joint-normalizer
polynomial identity of \cite{SandersGish2026FourCore} (Theorem~2),
where $(v, h, \beta, r)$ are the simplex coordinates on the
$4$-core $\mathcal{C} = \{0, 7, 8, 9\}$. Evaluation of the
unnormalized polynomial $Z$ at the $\mathrm{HARMONY}$-pure point
$(v, h, \beta, r) = (0, \mathrm{HARMONY}, 0, 0)$ gives
\[
  Z(0, \mathrm{HARMONY}, 0, 0) \;=\; \mathrm{HARMONY}^{2}.
\]
Normalizing by the triadic substrate volume $|\Z/10|^{3} = 1000$
yields $\Omega_b = \mathrm{HARMONY}^{2}/|\Z/10|^{3} = 49/1000$.
\end{theorem}

% [BRAYDEN-DERIVE]: The following paragraph requires Brayden's
% structural notes on why the HARMONY-pure projection of the
% 4-core normalizer corresponds to the "baryonic anchor" of the
% cosmological substrate. The mathematical evaluation Z(0,7,0,0) = 49
% is trivial; the COSMOLOGICAL interpretation is the substantive
% content.

\paragraph{Substrate interpretation.}
[BRAYDEN-DERIVE: needs explicit structural reading of the
HARMONY-pure projection of the $4$-core as the cosmological
``baryonic anchor.'' Provisional claim, requiring derivation:
the $\mathrm{HARMONY}$-pure projection corresponds to the
substrate state with all mass concentrated on the
$\sigma$-fixed-yet-fully-coupled vertex of the $4$-core, which
in the BEING/DOING/BECOMING triadic structure is the
``stable anchor'' that physical baryonic matter inherits via
the discrete-to-continuum reduction. The resulting count
$\mathrm{HARMONY}^{2} = 49$ is the squared $\mathrm{HARMONY}$
mass at this projection point.]

\subsection{$\Omega_{DM}$ from the becoming-layer count}

\begin{theorem}[$\Omega_{DM}$ derivation]\label{thm:omega-dm-deriv}
Let $V$ denote the $4$-dimensional commutative non-associative
$\F_5$-algebra obtained by taking the $\F_5$-bilinear extension of
the $\mathrm{CL}$ fuse table restricted to the $4$-core
$\mathcal{C} = \{0, 7, 8, 9\} \subset \Z/10$ and reducing mod $5$
to representatives $\{0, 2, 3, 4\} \subset \F_5$. Then
\[
  |V| = \dim_{\F_5} V = 4,\qquad |\Aut(V)| = 40,
\]
where $|\Aut(V)| = 40$ is verified by direct constraint enumeration
(\texttt{verify\_aut\_V\_order.py}) and identified group-theoretically
in \cite{SandersClaudeChat2026BridgeSprint} (WP118) as the order of
$F_{20} \times \Z/2$, where $F_{20}$ is the Frobenius group of
order $20$. With $|\sigma| = 6$ the $\sigma$-cycle length on
$\Z/10 \setminus \{0, 3, 8, 9\}$ from
\cite{SandersGishJohnson2026SigmaRate} (Theorem~A),
\[
  \Omega_{DM}
  \;=\; \frac{(|\Aut(V)| + |V|) \cdot |\sigma|}{|\Z/10|^{3}}
  \;=\; \frac{(40 + 4) \cdot 6}{1000}
  \;=\; \frac{264}{1000}.
\]
\end{theorem}

% [BRAYDEN-DERIVE]: The numerator (40 + 4) * 6 = 264 admits a
% structural reading as "becoming-layer harmonic count times
% sigma-cycle length", but this reading needs explicit derivation
% from the substrate's triadic structure. Specifically:
%   - Why is the becoming-layer count |Aut(V)| + |V| = 44?
%   - Why does the sigma-cycle factor in linearly (rather than
%     squared, cubed, or as a permutation orbit count)?
%   - Why does the becoming-layer (rather than the being-layer
%     or doing-layer) carry the cosmological dark-matter density?

\paragraph{Substrate interpretation.}
[BRAYDEN-DERIVE: needs explicit structural reading of the
becoming-layer count $|\Aut(V)| + |V| = 44$ and the linear
$\sigma$-cycle factor. Provisional claim: the becoming-layer
of the BEING/DOING/BECOMING triadic structure carries the
$4$-core's automorphism orbits ($|\Aut(V)|$ of them) plus its
$|V|$ basis vectors, totaling $44$ symmetry-distinguished cells;
the $\sigma$-cycle factor counts the trajectories of each cell
under the substrate's natural permutation. The cosmological
dark-matter density is the becoming-layer's symmetry-orbit count
relative to the triadic substrate volume.]

\subsection{$\Omega_\Lambda$ by closure}

\begin{theorem}[$\Omega_\Lambda$ derivation]\label{thm:omega-l-deriv}
Closure $\Omega_b + \Omega_{DM} + \Omega_\Lambda = 1$ together
with Theorems~\ref{thm:omega-b-deriv} and \ref{thm:omega-dm-deriv}
gives
\[
  \Omega_\Lambda
  \;=\; 1 - \Omega_b - \Omega_{DM}
  \;=\; \frac{|\Z/10|^{3} - \mathrm{HARMONY}^{2} - (|\Aut(V)| + |V|)|\sigma|}{|\Z/10|^{3}}
  \;=\; \frac{687}{1000}.
\]
\end{theorem}

\paragraph{Cubic-anchor reading.}
The numerator $687 = 2 \cdot 343 + 1 = 2 \cdot \mathrm{HARMONY}^{3} + 1$
admits a structural reading as ``the cubic $\mathrm{HARMONY}$
anchor doubled, plus a closure-completing offset of unity.''
The doubling reflects matter--antimatter pair structure of the
substrate; the $+1$ offset is the structural-completeness
constant required by exact triadic closure to $|\Z/10|^{3}$.
Both readings are consistent with the closure derivation above
and emphasize different facets of the same integer identity.

%==============================================================
\section{The dark-energy scale $\Lambda \approx 1.74$ meV}\label{sec:lambda}
%==============================================================

In the freezing-quintessence model of the JCAP companion
paper~\cite{SandersGishJohnson2026JCAP}, the action is
\[
  S = \int d^{4}x\,\sqrt{-g}\Big[\frac{R}{16\pi G}
     - \frac{M_\mathrm{Pl}^{2}}{2}\,g^{\mu\nu}\partial_\mu\Xi\,\partial_\nu\Xi
     - \Lambda^{4}\,\Xi\log\Xi\Big]
\]
with mass scale $\Lambda$, dimensionless field $\Xi$, and analytic
vacuum at $\Xi_0 = e^{-1}$ giving fluctuation mass
$m_\Xi^{2} = \Lambda^{4} e / M_\mathrm{Pl}^{2}$.

\begin{theorem}[$\Lambda$-scale relation]\label{thm:lambda-scale}
With Friedmann normalization $\rho_{c,0} = 3 H_{0}^{2} M_\mathrm{Pl}^{2}$
and the dark-sector trinity (Theorem~\ref{thm:trinity}),
\[
  \frac{\Lambda^{4}}{\rho_{c,0}}
  \;\stackrel{?}{=}\;
  \frac{\Omega_\Lambda}{3}
  \;=\;
  \frac{2\cdot\mathrm{HARMONY}^{3} + 1}{3 \cdot |\Z/10|^{3}}
  \;=\; \frac{687}{3000}
  \;=\; 0.229,
\]
giving $\Lambda \approx 1.74$~meV. The fitted value in
\cite{SandersGishJohnson2026JCAP} is $\Lambda \approx 1.7$~meV,
matching the structural prediction at $2.5\%$.
\end{theorem}

% [BRAYDEN-DERIVE]: The factor of 1/3 in the denominator is the
% standard Friedmann normalization rho_c0 = 3 H_0^2 M_Pl^2, but
% the relation rho_DE,0 = 3 Lambda^4 implicit in this identification
% is NOT a direct consequence of the action -- it depends on the
% present field configuration (Xi_today, Xi_dot_today). A
% first-principles derivation of this relation requires either:
%   - the triadic-projection reading: cosmological observables are
%     one of three substrate layers (BEING/DOING/BECOMING), and the
%     1/3 factor is the layer-projection coefficient; OR
%   - the freezing-branch energy budget: at the freezing trajectory
%     selected by the model, V_eff(Xi_today) = Lambda^4/3 to leading
%     order in (1 - Xi_today/Xi_0); OR
%   - some other identifiable structural relation.
% This is the central open question of the present paper.

\paragraph{The $1/3$ factor and the open question.}
The relation $\rho_\mathrm{DE,0} = 3 \Lambda^{4}$ implicit in
Theorem~\ref{thm:lambda-scale} is not a direct consequence of
the action; it depends on the present field configuration
$(\Xi_\mathrm{today}, \dot\Xi_\mathrm{today})$ and the Friedmann
history. We mark the question with $\stackrel{?}{=}$ in
Theorem~\ref{thm:lambda-scale} to be explicit. Empirically,
the JCAP companion fit~\cite{SandersGishJohnson2026JCAP}
reports $\Lambda^{4}/\rho_{c,0} = 0.231$, which together with
the Planck~2018 value $\Omega_\Lambda = 0.6862$ implies
\[
  \frac{\rho_\mathrm{DE,0}}{\Lambda^{4}}
  \;=\; \frac{\Omega_\Lambda \cdot \rho_{c,0}}{\Lambda^{4}}
  \;=\; \frac{0.6862}{0.231}
  \;=\; 2.971,
\]
within $1\%$ of $3$. The structural prediction of
Theorem~\ref{thm:lambda-scale}, $\rho_\mathrm{DE,0} = 3 \Lambda^{4}$
exactly, is consistent with the empirical late-time
configuration of the JCAP model at this precision. Three
candidate derivations of the $1/3$ factor are:

\begin{enumerate}
\item \textbf{Triadic projection.} Cosmological observables
emerge from one of three layers of the substrate's triadic
structure (BEING $\times$ DOING $\times$ BECOMING). The $1/3$
factor in $\Lambda^{4}/\rho_{c,0}$ is the projection coefficient
from the three-layer substrate to the one-layer cosmological
observable. This reading aligns with the $|\Z/10|^{3}$ denominator
in $\Omega_b$, $\Omega_{DM}$, $\Omega_\Lambda$ being the cube of
the substrate volume; the cosmological match recovers the cube
because Friedmann sums over all three layers, and the $1/3$
factor in $\Lambda^{4}$ recovers a single-layer projection.

\item \textbf{Freezing-branch energy budget.} On the freezing
trajectory identified in \cite{SandersGishJohnson2026JCAP}, the
present field configuration $(\Xi_\mathrm{today},
\dot\Xi_\mathrm{today})$ might satisfy
$V_\mathrm{eff}(\Xi_\mathrm{today}) \approx \Lambda^{4}/3$ to
leading order in $(1 - \Xi_\mathrm{today}/\Xi_{0})$. Verifying
this requires explicit FRW solution at the documented fit; the
manuscript presents this as a target for future work.

\item \textbf{Discrete-to-continuum bridge.} The $\Xi$ field of
\cite{SandersGishJohnson2026JCAP} carries the cosmological
$\Omega_\Lambda$ in the continuum limit; its discrete-substrate
origin (a normalized linear combination of the substrate's
$10$-operator state space) might fix the $1/3$ factor through
the explicit projection. This is the deepest of the three
candidates and is the open question of \S\ref{sec:open}.
\end{enumerate}

%==============================================================
\section{Open questions and future work}\label{sec:open}
%==============================================================

\subsection{The discrete-to-continuum bridge for $\Xi$}\label{sec:bb-bridge}

[BB-BRIDGE: needs explicit construction of the $\Xi$ scalar
field as a continuous projection of the substrate's $10$-operator
state space.]

The $\Xi$ field of \cite{SandersGishJohnson2026JCAP} carries
the cosmological $\Omega_\Lambda$ in the continuum limit, but its
discrete-substrate origin is not yet specified. Candidate
identifications include:
\begin{itemize}
\item $\Xi$ as the normalized projection onto a specific
$2$-dimensional subspace of the $10$-operator state space (e.g.,
the $\mathrm{BALANCE} \pm \mathrm{CHAOS}$ mode).
\item $\Xi$ as the trace of the substrate's density operator
along the $\sigma$-cycle.
\item $\Xi$ as the squared norm of the
$\mathrm{HARMONY}$-projected component of the substrate state.
\end{itemize}
Each candidate would predict a specific relation between the
discrete substrate's dynamics and the continuum field's potential
$V(\Xi) = \Lambda^{4} \Xi \log \Xi$. The Bialynicki-Birula--Mycielski
nonlinearity \cite{BialynickiBirulaMycielski1976}, which uniquely
characterizes $\Xi \log \Xi$ as the separability-preserving
nonlinearity in nonlinear quantum mechanics, suggests the
projection respects the substrate's tensor-product structure.

\subsection{An independent prediction: the scalar spectral index}\label{sec:prediction}

The dark-sector trinity (Theorem~\ref{thm:trinity}) and the
$\Lambda$-scale relation (Theorem~\ref{thm:lambda-scale}) are
constructed to match Planck~2018 dark-sector observations
$\Omega_b$, $\Omega_{DM}$, $\Omega_\Lambda$. To distinguish the
present framework from numerology we exhibit one observable that
was \emph{not} used in the construction yet falls within
$0.1\sigma$ of the published Planck value: the scalar spectral
index of primordial perturbations.

\begin{theorem}[Spectral-index prediction]\label{thm:ns}
With $\mathrm{HARMONY} = 7$ and $|\Z/10| = 10$ as in
\S\ref{sec:substrate}, the substrate predicts
\[
  n_s \;=\; 1 \;-\; \frac{\mathrm{HARMONY}}{2 \cdot |\Z/10|^{2}}
  \;=\; 1 \;-\; \frac{7}{200} \;=\; 0.9650.
\]
\end{theorem}

The Planck~2018 published value is $n_s = 0.9649 \pm 0.0042$
\cite{Planck2018}, in agreement with
Theorem~\ref{thm:ns} at $0.024\sigma$ ($\sim 0.01\%$ residual).
This prediction was first noted in the bridge-sprint working
material~\cite{SandersClaudeChat2026BridgeSprint} (WP125), where
it is identified as the small-amplitude tilt arising from the
substrate's $\mathrm{HARMONY}/|\Z/10|^{2}$-suppression of
quadrupole-mode mixing in an effective inflaton sector built on
the same $V$ algebra. The structural reading and the inflaton
construction are sketched in
\cite{SandersClaudeChat2026BridgeSprint} and remain open for
rigorous derivation in a forthcoming companion.

\paragraph{Why this counts as independent.}
The fit-input observables used to motivate the trinity are
$\Omega_b$, $\Omega_{DM}$, $\Omega_\Lambda$. The spectral index
$n_s$ is a separate, primordial CMB observable not entering
either the closure relation
$\Omega_b + \Omega_{DM} + \Omega_\Lambda = 1$ or the three
Hubble-independent ratio tests (\S\ref{sec:trinity}).
The same two integer primitives $\mathrm{HARMONY}$ and
$|\Z/10|$ that produce $\Omega_b$ and $\Omega_\Lambda$ also
produce $n_s$ at $0.024\sigma$, with no parameter freedom
between observables.

\paragraph{Other candidates.}
The same bridge-sprint material identifies further candidate
predictions in the same primitives:
\begin{itemize}
\item Baryon-photon ratio $\eta = |\sigma|/|\Z/10|^{10}$
predicting $\eta = 6 \times 10^{-10}$ vs measured
$\eta \approx 6.1 \times 10^{-10}$ ($1.6\%$).
\item The Cabibbo-angle refinement
$\lambda_C = (|\sigma| + |\mathcal{C}|)/(\mathrm{HARMONY}^2)$
$= 11/49 \approx 0.224$.
\item The matter-DE transition redshift $z_{eq}^{(\Lambda)}
\approx 0.30$ from $\Omega_{DM}/\Omega_\Lambda \cdot
\mathrm{HARMONY}/|\sigma|$.
\end{itemize}
We report these candidates as supporting the framework's
predictive reach but do not develop them in the present
submission, since each requires a separate physical-side
derivation analogous to that for $n_s$.

\subsection{$\Omega_{DM}$ derivation rigor}

[BRAYDEN-DERIVE: $\Omega_{DM}$ derivation in
\S\ref{sec:derivations} requires the explicit becoming-layer
construction. The numerical match $264 = 44 \cdot 6$ is
arithmetically trivial; the structural derivation is open.]

\subsection{Three-paper foundation triangle}

The present paper is the third of three companion submissions
that, together, encode the dark-sector decomposition into
algebraic ingredients. The $\sigma$-rate paper
\cite{SandersGishJohnson2026JCAP} establishes the substrate's
$|\Z/10|^{3}$ normalization. The $4$-core seed paper
\cite{SandersGish2026FourCore} establishes the
$(v + h + \beta + r)^{2}$ normalizer that gives
$\mathrm{HARMONY}^{2}$ on the $\mathrm{HARMONY}$-projection. The
present paper combines these structural pieces with an explicit
derivation of $\Omega_{DM}$ and the $\Lambda$-scale relation,
closing the trinity.

%==============================================================
\section{Verification}\label{sec:verify}
%==============================================================

The numerical match (Theorem~\ref{thm:trinity}, Table~\ref{tab:planck},
Theorem~\ref{thm:uniqueness}) is verified by two scripts at
DOI~\texttt{10.5281/zenodo.18852047}:

\begin{itemize}\setlength\itemsep{0pt}
\item \texttt{sprint18\_uniqueness\_search.py} performs the main
verification: direct evaluation of $\Omega_b$, $\Omega_{DM}$,
$\Omega_\Lambda$ at $(H, N) = (7, 10)$; comparison to Planck~2018
published values with explicit $\sigma$-residuals; the three
Hubble-independent ratio tests; the uniqueness search across $784$
small-integer $(H, N)$ pairs and the substrate-factorization check
for the six closure-exact candidates.

\item \texttt{verify\_aut\_V\_order.py} verifies the supplementary
fact $|\Aut(V)| = 40$ used in
Theorem~\ref{thm:omega-dm-deriv}: it reconstructs the $4$-core's
$\F_5$-lift $V$ from the CL fuse table directly, enumerates
all matrices preserving $V$'s multiplication, and confirms that
exactly $40$ such matrices exist.
\end{itemize}

Both scripts run to completion in under five seconds and print all
verification numbers cited in the manuscript.

%==============================================================
\begin{thebibliography}{99}
%==============================================================

\bibitem{Planck2018}
N.~Aghanim et al.\ (Planck Collaboration),
\emph{Planck 2018 results.\ VI.\ Cosmological parameters},
A\&A \textbf{641}, A6 (2020).

\bibitem{BialynickiBirulaMycielski1976}
I.~Bialynicki-Birula and J.~Mycielski,
\emph{Nonlinear wave mechanics},
Annals of Physics \textbf{100}(1--2), 62--93 (1976).

\bibitem{SandersGish2026FourCore}
B.~R.~Sanders and M.~Gish,
\emph{Joint Closure of Two Commutative Binary Operations on
$\mathbb{Z}/10\mathbb{Z}$: an Eight-Element Chain and a
Normalizer Identity},
Submitted to Communications in Algebra, 2026;
Zenodo DOI 10.5281/zenodo.18852047.

\bibitem{SandersGishJohnson2026JCAP}
B.~R.~Sanders, M.~Gish, and H.~J.~Johnson,
\emph{Logarithmic Quintessence: A Dimensionless Scalar Dark
Energy Model with an Analytic Vacuum},
In preparation, 2026;
Zenodo DOI 10.5281/zenodo.18852047.

\bibitem{SandersGishJohnson2026SigmaRate}
B.~R.~Sanders, M.~Gish, and H.~J.~Johnson,
\emph{Non-Associativity Decay in Binary Composition Tables over
$\mathbb{Z}/N\mathbb{Z}$},
Submitted to Journal of Combinatorial Theory, Series~A, 2026;
Zenodo DOI 10.5281/zenodo.18852047.

\bibitem{SandersClaudeChat2026BridgeSprint}
B.~R.~Sanders et al.,
\emph{Bridge Sprint: Discrete Dirac on the $4$-Core's
$\F_p$-Lift, F$_p$-Universality, and Dark-Sector Algebraic
Predictions},
Sprint working bundle (WP117--WP127), 2026-05-04;
Zenodo DOI 10.5281/zenodo.18852047.

\bibitem{Sprint18FrontierPredictions}
B.~R.~Sanders et al.,
\emph{[Future companion paper on independent
substrate-cosmology predictions]},
In preparation, 2026.

\end{thebibliography}

\end{document}
